\section{Discussion}
\label{sec:discussion}

{\color{red}%
\DIFdel{%
Limitations. First, the performance benefits vary across kernel types and hardware platforms. TilingInfer is most effective for lightweight kernels with high invocation frequency on Ascend, while on GPUs, further specialization on top of existing manual optimization yields limited improvement due to sufficient general-purpose registers and well-optimized baseline libraries.
}
\par}

{\color{blue}%
Workload-dependent benefits. Our experiments show that the magnitude of TilingInfer's benefit depends strongly on workload characteristics. Recommendation and decode workloads benefit most because Attention and MatMul expose profitable invariant schedules, yielding about 10\% average inference speedup. Schedule-fixed expert-written libraries also benefit from removing unreachable template instances. This per-model pruning is most attractive when a deployment serves fewer models, because each specialized binary can discard a larger fraction of the universal library. By contrast, HSTU and prefill remain compute-heavy, so specializing their core Attention and MatMul kernels removes little of the total work and improves end-to-end latency by only about 2\%. Some kernels in the current open-source baselines also cannot be further specialized by TilingInfer, further limiting end-to-end gains.
\par}

{\color{red}%
\DIFdel{%
Second, the current prototype is based on PyTorch; applying TilingInfer to other frameworks like TensorFlow requires constructing update rules for their specific data types.
}
\par}

{\color{blue}%
Platform applicability. Our schedule-generic kernel-specialization implementation currently targets Ascend. In the open-source libraries we surveyed for other accelerators, expert-written kernels predominantly follow the schedule-fixed design, leaving no comparable schedule-generic baseline to evaluate. Ascend exposes many programmable hardware levels; achieving high performance therefore relies on numerous runtime schedule parameters, which motivates schedule-generic libraries and makes automatic specialization particularly relevant.
\par}

{\color{red}%
\DIFdel{%
Third, while TilingInfer can potentially specialize Host-side programs, such optimization involves more complex control flows and external function calls, which we leave for future work.
}
\par}

{\color{blue}%
Optimization scope and limitations. TilingInfer applies to existing expert-written operator libraries: it specializes schedule-generic kernels or removes unreachable instances from schedule-fixed libraries. This differs from AI compilers, which generate implementations and search schedules; their code-generation pipelines do not optimize an existing C++ expert library in place.
\par}

{\color{blue}%
Besides, TilingInfer analyzes host code only to recover constants and launch-site liveness, then optimizes reachable device kernels; it does not rewrite host code. Host control flow is complex, CPUs already execute it efficiently, and profitable host specialization would require a more selective policy than device-kernel constant propagation. Production NPU Graph execution further suppresses CPU overhead. In layered Transformers, one host path is typically executed once per decode step while its device kernels repeat across many layers, making kernel optimization the higher-value target.
\par}
